\section{Iterated protection closure and a strict original far-head surplus}
\label{sec:iterated-protection-closure-strict-surplus}

All graphs are finite, nonempty, loopless and without directed two-cycles.
First and second neighborhoods always mean directed distance exactly one
and exactly two. For any graph $H$, protection means
\[
 a\triangleleft_H w \quad\Longleftrightarrow\quad
 a\to w\in A(H),\qquad
 N_H^-(a)\subseteq N_H^-(w)\setminus\{a\}.
\]
The protected vertex $a$ is far from $w$. A missing direction $a\to b$
is forced in $H$ when some $w$ satisfies $a\triangleleft_H w$ and
$w\to b\in A(H)$. As previously proved, a forced direction is uniquely
convenient in $H$: every predecessor of $a$ already reaches $b$ within
two steps of $H$.

\subsection{A finite closure which preserves all protections already obtained}

Starting at $Q_0=D$, add \emph{all} directions currently forced in $Q_i$
to form $Q_{i+1}$. Stop when none remain, and call the final graph $Q$.
No arc is removed. Every nonterminal round fills at least one missing
pair, so there are at most $M(D)$ nonterminal rounds. The construction is
deterministic and never attempts to add two opposite directions.

\begin{lemma}[Monotone protection and a stable partial median]
\label{lem:iterated-closure-monotone-protection}
Every protection of $Q_i$ survives in $Q_{i+1}$. Consequently every added
arc has a tail which protects some vertex in the final graph $Q$.
In $Q$, every protected pair has both incoming and outgoing neighborhood
nesting. Every global median order of $Q$ therefore extends its protection
partial order.
\end{lemma}
\begin{proof}
Consistency and convenience hold in each current graph by the forced
direction lemma; no minimality or deficiency is used.
Fix $y\triangleleft_{Q_i}x$ and a newly incoming arc $a\to y$.
There is a witness $a\triangleleft_{Q_i}w$, $w\to y$ in $Q_i$.
Protection of $y$ implies $w\to x$ in $Q_i$. If $ax$ is missing,
the same witness forces $a\to x$ in this round. An old $x\to a$
would imply $x\to w$, contradicting $w\to x$. Thus the new incoming
arc cannot destroy the required nesting at $y,x$.

At termination, if $y\triangleleft_Q x$ and $x\to z$ in $Q$, then
$z\to y$ would contradict incoming nesting, whereas a missing $yz$
would still force $y\to z$. Therefore $y\to z$ is already in $Q$.
This proves outgoing nesting. Swapping a reversed pair $x,y$ in a
vertex order gains the mutual arc $y\to x$, and loses no intermediate
incoming or outgoing contribution. Missing arcs contribute zero, so
this argument applies to the incomplete graph $Q$. Every global
median order of $Q$ places $y$ before $x$.
\end{proof}

\subsection{Maximization after closure, not during its construction}

Orient every remaining missing pair of $Q$ forward in one linear
extension of its protection order. Since there are no forced pairs left,
the previously proved protection-preserving reference theorem applies
with an empty forced set. Each resulting tournament preserves every
protection of $Q$, hence every protection of every earlier $Q_i$.
Call the still-missing pairs \emph{free pairs} in this subsection.

Choose a tournament $T_*$ maximizing its ordinary global median score
over \emph{all} these reference extensions. No assumption is made that
an intermediate $Q_i$ is all-deficient or arc-minimal.

\begin{lemma}[All-median rigidity and exact original feed neighborhoods]
\label{lem:iterated-closure-original-feed-neighborhoods}
Every global median order of $T_*$ extends protection in $Q$ and places
every free added arc forward. If $f$ is any of its median feeds, then
\[
 N_{T_*}^+(f)=N_Q^+(f)=N_D^+(f),\qquad
 N_Q^{++}(f)=N_D^{++}(f).
\]
Moreover, writing $r$ for the number of free pairs,
\[
 \max_T\operatorname{medscore}(T)
   =\operatorname{medscore}(Q)+r.
\]
The possible feeds across all maximizing completions are precisely the
global median feeds of the fixed partial graph $Q$.
\end{lemma}
\begin{proof}
Tournament protection gives domination, so the two-position swap forces
every median order to extend protection. If a free added arc were
backward in any median of a maximizing $T_*$, using that median as a new
reference would flip all backward free arcs forward without changing any
$Q$ arc. This strictly increases the score, a contradiction. Closure
under \emph{all} protection extensions is essential here.

A feed cannot have an outgoing free arc because it would be backward.
Nor can it be the tail of any arc added during closure: that tail has
a protection successor which survives in $Q$ and $T_*$, and must occur
later in every median order. Thus no added first arc exists at $f$,
at any stage of the construction.

Induct on the closure stages to obtain exact second-neighborhood equality.
Any new path from $f$ uses an original first arc $f\to a$. If its
second arc $a\to b$ was added in the current round, its convenience in
the \emph{previous} graph means that this predecessor $f$ already reached
$b$ within two steps there. A first target remains first throughout, so
no new exact second head occurs. Old exact second paths persist because
the first neighborhood is unchanged. The induction proves the equality
with $D$. It does not assert that every closure arc is convenient in $D$.

Each $Q$ median extends protection, and completing all free pairs forward
adds exactly $r$ to its score. Conversely any tournament/order pair scores
at most $\operatorname{medscore}(Q)+r$. Equality requires both components
to attain their maxima. This also gives the feed equivalence.
\end{proof}

\subsection{An original positive-gap feed cannot exhaust its far budget}

Let $\mathcal F_D(f)$ denote the originally far heads not protected in
$D$, let $t_D(f)=|\mathcal F_D(f)|$, and let $\mathcal M_Q(f)$ denote
the missing partners of $f$ in the terminal partial graph $Q$.
Put $g_H(f)=|N_H^+(f)|-|N_H^{++}(f)|$ and
\[
 A_f=N_{T_*}^{++}(f)\setminus N_D^{++}(f).
\]
The ordinary strong feed theorem and its global equality sedimentation
are those of Havet--Thomasse; no arc-weighted feed theorem is used.

\begin{theorem}[Strict surplus, without a minimal-counterexample hypothesis]
\label{thm:iterated-closure-strict-original-surplus}
For any median feed $f$ of $T_*$ with $g_D(f)>0$, one has
\[
 \mathcal M_Q(f)\ne\varnothing,\qquad
 t_D(f)\ge g_D(f)+1.
\]
More precisely, if $g_{T_*}(f)=0$, then
\[
 \varnothing\ne\mathcal M_Q(f)
     \subseteq U_{T_*}(f)\cap(\mathcal F_D(f)\setminus A_f),
 \qquad
 t_D(f)\ge g_D(f)+|\mathcal M_Q(f)|.
\]
This is an arbitrary-graph theorem: it needs neither all-deficiency of
$D$, strong connectivity, residual assumptions, nor arc minimality.
\end{theorem}
\begin{proof}
Protection in $D$ survives. The exact first set is original and original
exact second paths persist, so
\[
 N_{T_*}^{++}(f)=N_D^{++}(f)\,\dot\cup\,A_f,
 \qquad A_f\subseteq\mathcal F_D(f),\qquad
 g_{T_*}(f)=g_D(f)-|A_f|\le0.
\]
The last inequality is the tournament feed theorem.

If $\mathcal M_Q(f)$ were empty, $f$ would be whole in $Q$. Every far
head $y$ of a whole vertex in an oriented graph is protected: $y\to f$;
if $z\to y$ then $f\to z$ would give $f\to z\to y$, so adjacency
at $f$ forces $z\to f$. Thus $y\triangleleft_Q f$. Since all these
protections survive in $T_*$, $f$ acquires no further second heads.
The original first and second sets then give $g_{T_*}(f)=g_D(f)>0$,
a contradiction. Therefore $\mathcal M_Q(f)$ is nonempty.

Suppose $g_{T_*}(f)=0$. For any median order $L$ ending at $f$, the
strong feed inequalities imply $G_L=N_{T_*}^{++}(f)$ and
$|G_L|=|N_{T_*}^+(f)|$. If a remaining missing partner $b$ belonged
to this second set, the completed direction would be the free incoming
arc $b\to f$, since $f$ has no added first arc. Global sedimentation
puts $f$ before every vertex of $G_L$ while retaining a global median
order of the \emph{same} $T_*$. That free arc becomes backward,
contradicting all-median rigidity.

Every partner in $\mathcal M_Q(f)$ is consequently far in $T_*$.
It was also an original missing partner, hence cannot have been an
original protected head, and belongs to $\mathcal F_D(f)\setminus A_f$.
Also $|A_f|=g_D(f)$ in the equality case. The two disjoint subsets
$A_f$ and $\mathcal M_Q(f)$ give the sharper budget inequality.

If $g_{T_*}(f)<0$, integrality gives
$|A_f|\ge g_D(f)+1$, which also proves the strict-surplus bound.
\end{proof}

\begin{corollary}[An unbounded sufficient structural condition]
\label{cor:iterated-closure-low-far-budget}
If every deficient vertex $v$ of an oriented graph $D$ satisfies
$t_D(v)\le g_D(v)$, then $D$ has a Seymour vertex.
\end{corollary}
\begin{proof}
Otherwise all original gaps are positive and any feed of the construction
contradicts the strict-surplus theorem.
\end{proof}

\begin{corollary}[Nonunit residual at every closure feed]
\label{cor:iterated-closure-nonunit-feed}
If $D$ belongs to the previously strongly connected arbitrary-
arc-set-deletion-minimal all-deficient class, every median feed of $T_*$
has $\eta\ge2$. If $t_D(f)\le2$ at such a feed, necessarily
$g_D(f)=1$ and $t_D(f)=2$.
\end{corollary}
\begin{proof}
The zero and unit residual budgets are $t=0$ and $t\le1$ respectively,
while the strict-surplus theorem gives $t\ge g_D+1\ge2$. For $t\le2$
the same inequalities force $g_D=1,t=2$.
\end{proof}

\paragraph{Comparison and scope.}
The two balanced branches of the preceding \emph{one-round} forced-family
analysis are now absent for the \emph{iterated-closure} maximizing family.
This does not retroactively claim that they are impossible in every old
one-round completion. The new terminal graph need not remain a minimal
counterexample, and its arcs need not all be convenient relative to $D$.
No minimality-based packet bound is applied to an intermediate graph.
The strict-surplus theorem still permits arbitrarily large original
unprotected far budgets. It neither proves the complete SNC nor supplies
an actual all-deficient graph. External novelty and peer review remain
outstanding.
